\medskip
	On considère la suite $(u_{n})$ définie par $u_{0}=5$ et, pour tout $n$ :
$u_{n+1}=2+\ln \left(u_{n}^{2}-3\right)$



	\subsection*{Partie A : Exploitation de programmes Python}

	\begin{enumerate}
		\item On complète le script Python ci-dessous pour que \texttt{suite(k)}  renvoie la liste des $k$ premières valeurs de la suite $(u_{n})$.





		\begin{center}
		\begin{ttfamily}
		\begin{tabularx}{0.6\linewidth}{|X|}\hline
			def suite(k):\\
			\quad	L = []\\
			\quad	u = 5\\
			\quad	for i in range({\blue k}):\\
			\qquad	L.append(u)\\
			\qquad	u = \blue 2 + log( u**2 - 3\\
			\quad	return({\blue L})\\ \hline
		\end{tabularx}
		\end{ttfamily}
		\end{center}
		
		\item On a exécuté \texttt{suite(9)} ci-dessous. 
		


\begin{center}
\begin{ttfamily}
		
\begin{tabularx}{\linewidth}{|X|}\hline
>{}>{}>{} suite(9) \\
$\left[\right.$ 5, 5.091042453358316, 5.131953749864703,\\
5.150037910978289, 5.157974010229213, 5.1614456706362954,\\
5.162962248594583, 5.163624356938671, 5.163913344065642$\left.\right]$\\ \hline

\end{tabularx}
\end{ttfamily}
\end{center}

La suite $(u_n)$ semble croissante, mais on ne peut pas se prononcer sur sa limite.

		\item On a ensuite créé la fonction \texttt{mystere(n)} donnée ci-dessous et exécuté \linebreak \texttt{mystere(10000)}, ce qui a renvoyé \texttt{1}.



\begin{center}
		\begin{ttfamily}

\begin{tabularx}{\linewidth}{|X|}\hline
			def mystere(n):\\
			\quad L = suite(n)\\
			\quad	c = 1\\
			\quad	for i in range(n - 1):\\
			\qquad	if L[i] $>$ L[i + 1]:\\
			\qquad \quad c = 0\\
			\quad return c\\ \hline

\end{tabularx}
		\end{ttfamily}
		\end{center}

La fonction \texttt{mystere} renvoie 0 s'il existe un indice $i$ pour lequel \texttt{L[i] > L[i+1]}, c'est-à-dire $u_{n}>u_{n+1}$. Comme \texttt{mystere(10000)} renvoie 1, cela veut dire que pour tout $i\leqslant \np{9999}$, on a: $u_{i} \leqslant u_{i+1}$.

Cela ne contredit pas la conjecture faite sur la croissance de la suite $(u_n)$.












	\end{enumerate}

\subsection*{Partie  B : Étude de la convergence de la suite ( $\boldsymbol{u}_{\boldsymbol{n}}$ )}

On considère la fonction $g$ définie sur $[2 ~;~+\infty [$ par :
$g(x)=2+\ln \left(x^{2}-3\right)$.



\begin{enumerate}
	\item 
$g'(x)= 0+ \dfrac{2x}{x^2-3} = \dfrac{2x}{x^2-3}$

Sur  $[2 ~;~+\infty[$, $2x>0$ et $x^2>4$ donc $x^2-3>1>0$;
donc $g'(x)>$ sur 	$[2 ~;~+\infty[$, ce qui prouve que la fonction $g$ est croissante sur $[2 ~;~+\infty[$.
	
	\item 
		\begin{enumerate}
		\item On démontre par récurrence que la propriété $4 \leqslant u_{n} \leqslant u_{n+1} \leqslant 6$ est vraie pour tout  entier naturel $n$.

\begin{list}{\textbullet}{}
\item \textbf{Initialisation}

$u_0=5$ et $u_1=2+\ln\left (5^2-3\right ) = 2+\ln(22) \approx 5,09$

Donc $4 \leqslant u_{0} \leqslant u_{1} \leqslant 6$; la propriété est vraie au rang 0.
\item \textbf{Hérédité}

On suppose la propriété vraie au rang $n$, c'est-à-dire $4 \leqslant u_{n} \leqslant u_{n+1} \leqslant 6$.

D'après la définition de la fonction $g$, on a $u_{n+1}=g(u_n)$ pour tout $n$.

On sait que la fonction $g$ est croissante sur $[2\;;\;+\infty[$; donc de l'inégalité  $4 \leqslant u_{n} \leqslant u_{n+1} \leqslant 6$, on déduit $g(4) \leqslant g\left (u_{n}\right ) \leqslant g\left (u_{n+1}\right ) \leqslant g(6)$.

$g(4)\approx 4,56$; $g\left (u_n\right )=u_{n+1}$; $g\left (u_{n+1}\right )=u_{n+2}$
et $g(6)\approx 5,5$

On a donc $4,5 \leqslant u_{n+1} \leqslant u_{n+2} \leqslant 5,5$,
ce qui entraine $4 \leqslant u_{n+1} \leqslant u_{n+2} \leqslant 6$.

La propriété est donc héréditaire.
\item \textbf{Conclusion}

La propriété est vraie au rang 0 et elle est héréditaire pour tout $n\geqslant 0$ donc, d'après le principe de récurrence, elle est vraie pour tout entier naturel $n$.
\end{list}

On a donc  démontré  que, pour tout $n$ : $4 \leqslant u_{n} \leqslant u_{n+1} \leqslant 6$.

		\item 
\begin{list}{\textbullet}{D'après la question précédente, pour tout $n$, on a:}
\item $u_{n}\leqslant u_{n+1}$ donc la suite $(u_n)$ est croissante;
\item $u_n \leqslant 6$ donc la suite $(u_n)$ est majorée.
\end{list}		

La suite $(u_n)$ est croissante et majorée donc, d'après le théorème de la convergence monotone, elle est convergente.
		
		
		\end{enumerate}
	\end{enumerate}

\subsection*{Partie C : Étude de la valeur de la limite}

On considère la fonction $f$ définie sur [ 2 ; $+\infty$ [ par :
$f(x)=2+\ln \left(x^{2}-3\right)-x$.



On donne le tableau de variations de $f$ suivant. 

\[\begin{tablvar}[6em]{2}
\hline
x & 2 && 3 && +\infty\\
\hline
\variations{\mil{f(x)} & \bas{0} && \haut{\ln(6)-1} &&  \bas{-\infty}} 
\hline
\end{tablvar}\]

	\begin{enumerate}
		\item \begin{enumerate}
			\item 
D'après le tableau de variations de $f$, on voit que $f(2)=0$. \\
On complète le tableau:

\[\begin{tablvar}[6em]{2}
\hline
x & \vr{2} && 3 & \vr{\beta} & +\infty\\
\hline
\variations{\mil{f(x)} & \posvr{3}{0} && \haut{\ln(6)-1} & \vr{0}&  \bas{-\infty}} 
\hline
\end{tablvar}\]

Donc l'équation $f(x)=0$ admet exactement deux solutions sur l'intervalle $[2 ~;~+\infty[$ que l'on notera $\alpha=2$ et $\beta \in \left [ 3\;; +\infty \strut\right [$

			\item 
On a vu que $\alpha=2$.

\begin{list}{\textbullet}{À la calculatrice, on trouve successivement:}
\item $f(5)<0<f(6)$ donc $\beta \in \left [ 5\;; 6 \strut\right ]$;
\item $f(5,1)<0<f(5;2)$ donc $\beta \in \left [ 5,1\;; 5,2 \strut\right ]$;
\item $f(5,16)<0<f(5;17)$ donc $\beta \in \left [ 5,16\;; 5,17 \strut\right ]$;
\item $f(5,164)<0<f(5;165)$ donc $\beta \in \left [ 5,164\;; 5,165 \strut\right ]$.
\end{list}			

Donc $5,164$  une valeur approchée à $10^{-3}$ près de $\beta$.
		\end{enumerate}

		\item On note $\ell$ la limite de la suite $(u_{n})$.


$f(x)=g(x) - x$ donc $g(x)=x \iff f(x)=0$

Pour tout $n$, on a $g\left (u_n\right )=u_{n+1}$ donc, par passage à la limite on a $g(\ell)=\ell$, ce qui équivaut à $f(\ell) = 0$.

$u_0 = 5$ et la suite $(u_n)$ est croissante donc $\ell\geqslant 5$.

$\ell\geqslant 5$ et $\ell$ vérifie $f(\ell) = 0$; d'après les questions précédentes, on peut en déduire que $\ell = \beta$.

	\end{enumerate}

\bigskip
